% Course: 77972 -- Physical Models of Biological Systems
% Academic year: 2026--2027 | Semester: A
% Lecture number: 1 | Document chapter: 0 (course introduction)
% Lecture date: 2026-10-13
% Title: Introduction
% Lecturer: Dr. Amir Erez
% Notes prepared by: Eran Rehani
% Source: Lecture Transcripts/77972_2026-10-13_שעור_א_Captions_עברית_Hebrew.md
% Main conceptual discussion: 00:10:48--00:37:06; modelling cycle: 00:48:54--00:49:18.
% Supporting references: models-notes-2026.pdf, chapter 0, printed pp. 9--11;
% Philip Nelson, Physical Models of Living Systems, chapter 2.
\chapter{Introduction}\label{chap:introduction}
\noindent\textbf{From Lecture 1 \enspace|\enspace 13 October 2026}

Physical modelling connects biological mechanisms to observations. A useful model makes a question precise enough to calculate and test, with explicit assumptions and limitations.

\section{Why living systems need selective descriptions}\label{sec:intro-complexity}
Living systems interact, adapt, and alter their surroundings: photosynthesis changed the atmosphere and hence the conditions for later life.

They maintain organization by exchanging matter and energy with their surroundings. This respects the second law: a steady biological state need not be thermodynamic equilibrium.

\section{What a mathematical model contributes}\label{sec:intro-model-purpose}
A model specifies \textbf{variables}, \textbf{processes}, and their \textbf{relations}, then derives consequences not already inserted by hand.
\begin{itemize}
\item \textbf{Consistency:} equations make a verbal explanation precise and resolve competing effects quantitatively.
\item \textbf{Prediction:} unexpected observations expose missing processes or faulty assumptions; precise biological forecasts remain difficult.
\item \textbf{Inference:} observations constrain hidden rates or interactions when direct measurement is difficult or perturbs the living system.
\item \textbf{Intuition:} a simple model isolates the dominant mechanism and helps us anticipate responses to change.
\end{itemize}

\section{Choose the scale, then test the assumptions}\label{sec:intro-description}
Choose among molecules, cells, populations, communities, and ecosystems. Do space, fluctuations, or discrete individuals matter? Retain the relevant mechanism; add detail when it changes the conclusion.

Data and mechanisms complement each other. Ice-cream searches reveal annual periodicity without modelling its cause: predictive correlation need not explain a mechanism.

\textbf{The modelling cycle:} formulate $\to$ solve $\to$ compare with data $\to$ revise. Agreement supports usefulness, rather than proving every assumption. Chapter~\ref{chap:lecture01} applies this cycle to hidden viral turnover.

\subsection*{How to use these notes for revision}
State assumptions, explain each equation term, and interpret results biologically. Attempt analytical homework before consulting solutions; these problems prepare you for the test. Vary simulation parameters to build intuition.

\subsection*{Further Reading}
\begingroup\small
Amir Erez, course notes, Chapter~0, pp.~9--11; Philip Nelson, \emph{Physical Models of Living Systems}, Chapter~2, ``Physics and Biology.''
\endgroup
